\section{Valg av utviklingsmetodikk}
\label{sec:uvikling-modell}

Prosjektet startet under to sentrale forutsetninger: teamet hadde ingen tidligere erfaring med C\# eller .NET, og kravspesifikasjonen fra HDO definerte et tydelig omfang med faste leveranser. Dette dannet grunnlaget for valg av arbeidsmetode.

En sekvensiell utviklingsmodell, som fossefallsmodellen, ble vurdert, men tidlig forkastet. Uten kjennskap til .NET-økosystemet var det ikke mulig å estimere implementasjonskompleksitet presist ved oppstart. Flere sentrale krav knyttet til håndtering av videostrøm krevde iterative avklaringer med HDO etter hvert som domeneulikhetene mellom AntMedia Server og NHN Video ble tydelige. Dette ville vært utfordrende i en modell der krav fastsettes tidlig.

Prosjektet fulgte derfor en iterativ arbeidsform inspirert av Scrum \cite{scrumguide}, der elementer som sprintinndeling og iterativ planlegging ble benyttet uten full implementering av rammeverkets seremonier og roller. Arbeidet ble organisert i ukentlige sprinter med tydelige mål, og resulterte fortløpende i kjørbar kode og grunnlag for videre prioritering. Dette ga fleksibilitet til å håndtere usikkerhet og løpende avklaringer, samtidig som fremdriften ble opprettholdt.